
\documentclass[11pt,a4paper]{article}
\usepackage[T1]{fontenc}
\usepackage{lmodern}
\usepackage{amsmath,amssymb,amsfonts,mathtools}
\usepackage{graphicx,booktabs,longtable,array,tabularx,adjustbox}
\usepackage{geometry,microtype,fancyhdr,lastpage,etoolbox,enumitem}
\usepackage{newunicodechar,textcomp}
\usepackage{hyperref}
\hypersetup{hidelinks,pdfcreator={XeLaTeX / QM1 academic edition}}
\geometry{left=20mm,right=20mm,top=22mm,bottom=23mm}
\setcounter{secnumdepth}{2}
\setlength{\parindent}{0pt}
\setlength{\parskip}{4pt plus 1pt}
\setlength{\emergencystretch}{2em}
\setlength{\tabcolsep}{3pt}
\renewcommand{\arraystretch}{1.16}
\setlist{nosep,topsep=3pt,itemsep=2pt,leftmargin=7mm}
\AtBeginEnvironment{longtable}{\footnotesize\setlength{\tabcolsep}{2.2pt}}
\newcommand{\tightlist}{\setlength{\itemsep}{0pt}\setlength{\parskip}{0pt}}
\providecommand{\Cov}{\operatorname{Cov}}
\providecommand{\Var}{\operatorname{Var}}
\providecommand{\Corr}{\operatorname{Corr}}
\providecommand{\MAD}{\operatorname{MAD}}
\providecommand{\E}{\mathbb{E}}
\providecommand{\R}{\mathbb{R}}
\newunicodechar{—}{\textemdash}
\newunicodechar{–}{\textendash}
\newunicodechar{€}{\texteuro}
\newunicodechar{≤}{\ensuremath{\leq}}
\newunicodechar{≥}{\ensuremath{\geq}}
\newunicodechar{×}{\ensuremath{\times}}
\newunicodechar{÷}{\ensuremath{\div}}
\newunicodechar{Σ}{\ensuremath{\Sigma}}
\newunicodechar{σ}{\ensuremath{\sigma}}
\newunicodechar{μ}{\ensuremath{\mu}}
\newunicodechar{Δ}{\ensuremath{\Delta}}
\newunicodechar{∑}{\ensuremath{\sum}}
\newunicodechar{≈}{\ensuremath{\approx}}
\newunicodechar{√}{\ensuremath{\surd}}
\newunicodechar{≠}{\ensuremath{\ne}}
\newunicodechar{→}{\ensuremath{\to}}
\newunicodechar{∞}{\ensuremath{\infty}}
\newunicodechar{∈}{\ensuremath{\in}}
\newunicodechar{−}{\ensuremath{-}}
\newunicodechar{²}{\ensuremath{^{2}}}
\newunicodechar{³}{\ensuremath{^{3}}}
\newunicodechar{₀}{\ensuremath{_{0}}}
\newunicodechar{₁}{\ensuremath{_{1}}}
\newunicodechar{₂}{\ensuremath{_{2}}}
\newunicodechar{₃}{\ensuremath{_{3}}}
\newunicodechar{₄}{\ensuremath{_{4}}}
\newunicodechar{ₜ}{\ensuremath{_{t}}}
\newunicodechar{ₚ}{\ensuremath{_{p}}}
\newunicodechar{ₙ}{\ensuremath{_{n}}}
\newunicodechar{ₓ}{\ensuremath{_{x}}}
\newunicodechar{ᵢ}{\ensuremath{_{i}}}
\newunicodechar{ŷ}{\ensuremath{\hat y}}
\newunicodechar{Ŷ}{\ensuremath{\widehat Y}}
\newunicodechar{ₘ}{\ensuremath{_{m}}}
\newunicodechar{ₑ}{\ensuremath{_{e}}}
\newunicodechar{ₐ}{\ensuremath{_{a}}}
\newunicodechar{ₖ}{\ensuremath{_{k}}}
\newunicodechar{ₛ}{\ensuremath{_{s}}}
\newunicodechar{ₗ}{\ensuremath{_{l}}}
\newunicodechar{ₕ}{\ensuremath{_{h}}}
\newunicodechar{ₒ}{\ensuremath{_{o}}}
\newunicodechar{₊}{\ensuremath{_{+}}}
\newunicodechar{₋}{\ensuremath{_{-}}}
\newunicodechar{₌}{\ensuremath{_{=}}}
\newunicodechar{₍}{\ensuremath{_{(}}}
\newunicodechar{₎}{\ensuremath{_{)}}}
\newunicodechar{⁻}{\ensuremath{^{-}}}
\newunicodechar{⁺}{\ensuremath{^{+}}}
\newunicodechar{⁰}{\ensuremath{^{0}}}
\newunicodechar{¹}{\ensuremath{^{1}}}
\newunicodechar{⁴}{\ensuremath{^{4}}}
\newunicodechar{⁵}{\ensuremath{^{5}}}
\newunicodechar{⁶}{\ensuremath{^{6}}}
\newunicodechar{⁷}{\ensuremath{^{7}}}
\newunicodechar{⁸}{\ensuremath{^{8}}}
\newunicodechar{⁹}{\ensuremath{^{9}}}
\makeatletter
\renewcommand{\section}{\@startsection{section}{1}{0pt}{2.0ex plus .7ex minus .3ex}{1.0ex}{\large\bfseries}}
\renewcommand{\subsection}{\@startsection{subsection}{2}{0pt}{1.8ex plus .5ex minus .2ex}{.7ex}{\normalsize\bfseries}}
\renewcommand{\subsubsection}{\@startsection{subsubsection}{3}{0pt}{1.3ex plus .4ex minus .2ex}{.5ex}{\normalsize\bfseries}}
\makeatother
\pagestyle{fancy}
\fancyhf{}
\renewcommand{\headrulewidth}{0.35pt}
\renewcommand{\footrulewidth}{0pt}
\fancyhead[L]{\small\textsc{KEDGE Business School --- QM1}}
\fancyhead[R]{\small\textsc{2026--2027}}
\fancyfoot[L]{\small Badr ABOUFARASSE, PhD}
\fancyfoot[R]{\small\thepage\ /\ \pageref{LastPage}}
\setlength{\headheight}{14pt}
\setkeys{Gin}{keepaspectratio}
\begin{document}

\begin{center}
{\small\textsc{KEDGE Business School \quad EBP1 \quad Quantitative Methods 1}}\\[6pt]
{\Large\bfseries Fiche 7 — Cas intégrateurs}\\[6pt]
{\normalsize\bfseries Énoncés des exercices}\\[6pt]
{\small Badr ABOUFARASSE, PhD \quad --- \quad 2026--2027}
\end{center}
\vspace{2mm}\hrule\vspace{6mm}

\hypertarget{pruxe9paration-aux-uxe9valuations}{%
\section{Préparation aux évaluations}\label{pruxe9paration-aux-uxe9valuations}}

Résoudre des exercices de contrôle continu et de partiel en autonomie.

\emph{Préparation CC / partiel --- 70 min}

\hypertarget{cas-a-pruxe9paration-au-contruxf4le-continu}{%
\subsection{Cas A -- Préparation au contrôle continu}\label{cas-a-pruxe9paration-au-contruxf4le-continu}}

Une école interroge 200 étudiants sur le prix maximal qu'ils accepteraient de payer pour un service optionnel :

\begin{longtable}[]{@{}lllllll@{}}
\toprule\noalign{}
Prix & {\([5,10[\)} & {\([10,15[\)} & {\([15,20[\)} & {\([20,30[\)} & {\([30,40[\)} & {\([40,60[\)} \\
\midrule\noalign{}
\endhead
\bottomrule\noalign{}
\endlastfoot
{\(n_i\)} & 10 & 25 & 45 & 60 & 40 & 20 \\
\end{longtable}

\begin{enumerate}
\item
  Identifier population observée, unité statistique, caractère et nature du caractère.
\item
  Construire le tableau complet : amplitudes, densités, fréquences, ECC, ECD et centres.
\item
  Tracer l'histogramme et les courbes cumulatives.
\item
  Déterminer la classe modale, la médiane, {\(Q_1\)}, {\(Q_3\)} et la moyenne approchée.
\item
  Calculer la variance de population, l'écart-type et le coefficient de variation.
\item
  Déterminer le prix maximal tel qu'environ 65\% des étudiants puissent payer au moins ce prix.
\item
  On demande aussi « satisfait / non satisfait ». Préciser la nature de cette variable et le graphique adapté.
\end{enumerate}

\emph{Préparation CC / partiel --- 110 min}

\hypertarget{cas-b-examen-blanc-inuxe9dit-110-minutes}{%
\subsection{Cas B --- Examen blanc inédit --- 110 minutes}\label{cas-b-examen-blanc-inuxe9dit-110-minutes}}

Cas autonome : les données ci-dessous ne sont pas celles des exercices des chapitres 4 à 6. Présenter toutes les étapes et justifier les arrondis.

\hypertarget{partie-1-double-ruxe9gression-40-min}{%
\subsubsection{Partie 1 --- Double régression (40 min)}\label{partie-1-double-ruxe9gression-40-min}}

\begin{longtable}[]{@{}lllllll@{}}
\toprule\noalign{}
t & 1 & 2 & 3 & 4 & 5 & 6 \\
\midrule\noalign{}
\endhead
\bottomrule\noalign{}
\endlastfoot
x & 5 & 6 & 8 & 9 & 11 & 13 \\
y & 12 & 15 & 18 & 20 & 24 & 29 \\
\end{longtable}

1. Construire le tableau t, x, y, t², x², y², tx, xy ; calculer les deux corrélations et les droites de x sur t puis y sur x. 2. Prévoir x et y pour t=8 ; qualifier la prévision.

\hypertarget{partie-2-suxe9rie-trimestrielle-40-min}{%
\subsubsection{Partie 2 --- Série trimestrielle (40 min)}\label{partie-2-suxe9rie-trimestrielle-40-min}}

\begin{longtable}[]{@{}lllll@{}}
\toprule\noalign{}
Année / Year & Q1 & Q2 & Q3 & Q4 \\
\midrule\noalign{}
\endhead
\bottomrule\noalign{}
\endlastfoot
A & 72 & 88 & 80 & 110 \\
B & 78 & 96 & 88 & 121 \\
\end{longtable}

1. Tracer la série ; calculer les MA4 et les CMA4. 2. Avec la tendance fournie Tₜ=2,8t+78, calculer les huit ratios Yₜ/Tₜ. 3. Calculer les moyennes par trimestre, les normaliser (somme 4) et prévoir t=9 à 12.

\hypertarget{partie-3-indices-30-min}{%
\subsubsection{Partie 3 --- Indices (30 min)}\label{partie-3-indices-30-min}}

\begin{longtable}[]{@{}lllll@{}}
\toprule\noalign{}
Produit / Product & p₀ & q₀ & p₁ & q₁ \\
\midrule\noalign{}
\endhead
\bottomrule\noalign{}
\endlastfoot
A & 10 & 40 & 12 & 45 \\
B & 8 & 60 & 9 & 55 \\
C & 20 & 20 & 22 & 25 \\
\end{longtable}

Construire les quatre agrégats ; calculer les indices de prix et de quantité de Laspeyres, Paasche et Fisher, puis l\textquotesingle indice de valeur ; vérifier la propriété de Fisher.

\end{document}
